% !TeX root = ../../../main.tex
\subsection{Introduction}
\label{ch04:sec:intro}


In many causal inference applications where we aim to quantify how well a treatment works, estimating \emph{heterogeneous treatment effects} (HTEs) could be more useful than only estimating population‐level \emph{average treatment effects} (ATEs), building on the intuition that the same treatment can vary in effectiveness when given to different individuals. 
% Consider, for example, in different clinical contexts, patients with the same disease might benefit from different medications. 
In survival analysis with right-censored outcomes (common in clinical trials and electronic health records), estimating HTEs can be especially challenging. In addition to the standard difficulties of causal inference (unobserved counterfactuals, confounding), the analyst must account for censoring, where the event of interest is only observed for a subset of subjects.  These features complicate identification and estimation, yet they are central in high-stakes applications such as precision medicine and individualized policy-making~\citep{zhu2020targeted,chapfuwa2021enabling,curth2021survite}.

\ChFourRebuttal{Recent years have seen a growing set of causal survival methods~\citep{chapfuwa2021enabling,curth2021survite,ch04:cui2023estimating,bo2024ameta,noroozizadeh2025impact,xu2024estimating,meir2025heterogeneous}.}
% Recent years have seen a growing set of causal survival methods, including deep representation-learning approaches~\citep{chapfuwa2021enabling, curth2021survite}, Causal Survival Forests~\citep{cui2023estimating}, survival meta-learners~\citep{bo2024ameta, noroozizadeh2025impact}, and outcome imputation methods~\citep{xu2024estimating, meir2025heterogeneous}. 
Despite methodological advancement, no standardized benchmark exists, limiting reproducibility and fair comparisons. 
Most studies rely on bespoke simulations or limited real datasets with unknown ground truth, with differing levels of censoring, survival distributions, and causal assumptions. As a result, comparisons are not standardized, the robustness of different proposed methods is unclear, and progress is difficult to measure. 
% Moreover, in part because of how relatively new causal survival methods are, papers that propose new such methods have not compared their proposed approaches against a comprehensive assortment of other causal survival methods.
% Separately, recent falsification work by~\citet{pmlr-v238-demirel24a} shows that applying survival HTE estimators to observational data without validating assumptions can be dangerously misleading, underscoring the need for a rigorously designed testbed.

\ChFourRebuttal{
% HTE benchmark (non-survival): \cite{NEURIPS2022_benchmark,Shimoni2018Benchmarking,Kapkic2024causalbench}; ATE survival causal benchmark: \cite{voinot2025causal}.
While there is a growing benchmarking literature for treatment-effect heterogeneity in fully observed outcomes (e.g., \cite{NEURIPS2022_benchmark,Shimoni2018Benchmarking,Kapkic2024causalbench}) and recent benchmarks for survival ATE estimation (e.g., \cite{voinot2025causal}), to our knowledge, there is not yet any benchmark for survival HTE estimation under right-censoring. This missing piece motivates our focus on heterogeneous effects in censored time-to-event data.
}


We introduce \textsc{SurvHTE‐Bench}, the first comprehensive benchmark for HTE estimation in right-censored survival data. Our contributions are as follows:
\begin{itemize}
    \item \textbf{Method unification:} We categorize existing survival HTE methods (and natural extensions of such existing methods that technically have not previously been published) into three broad families: outcome imputation methods, direct-survival causal methods, and survival meta-learners. We provide a modular implementation of 53 methods among these families. This is the first systematic framework that unifies survival HTE methods, facilitating reproducibility and extensibility.  
    % No existing causal survival analysis paper has experimented on as large a collection of causal survival estimators as what we present. Note that some of the causal survival estimators we present have not actually been published before but are common-sense extensions of existing approaches.
    \item \textbf{Synthetic benchmark design:} We present a curated suite of 40 synthetic datasets spanning eight causal configurations (with different combinations of randomization, unobserved confounding, overlap violation, informative censoring) crossed with five survival scenarios (with different survival and censoring distributions), yielding controlled settings with known ground-truth HTEs under realistic assumption violations.
    \item \textbf{Semi-synthetic and real data:} We also include 10 semi-synthetic datasets adapted from existing literature (real covariates with simulated treatments and outcomes) that aim to be more realistic compared to purely synthetic datasets while still having ground truth on HTEs. We further include 2 widely studied real datasets: the Twins dataset that has known ground truth \citep{almond2005the} (i.e., per twin, one has the treatment and the other does not, so that we observe both counterfactual outcomes), and the HIV clinical trial dataset without known ground truth \citep{hammer1996atrial}.
    % \item We provide systematic comparisons of when each estimator excels or fails, providing practical guidance and highlighting open methodological gaps. \textcolor{red}{add some key takeaway from experiments.} We find S-Learner in the survival meta-learners family demonstrates remarkable stability across challenging scenarios involving assumption violations and high censoring rates. \textcolor{red}{(update this with synthetic data result: choice of imputation methods is important for outcome imputation causal survival methods? Causal Survival Forests shows stable and competitive performance without extensive tuning on base model options, providing reliable "out-of-the-box" options for practitioners.}
    \item \textbf{Comprehensive evaluation:} We compare representative estimators across all settings. Our results show that no single method dominates: performance depends on causal assumptions, censoring, and survival dynamics. Notably, \textcolor{black}{S- and matching-learners} among survival meta-learners demonstrate robustness under severe violations and high censoring.
    % Our experiments reveal that the S-Learner and \textcolor{red}{Matching} method in the survival meta-learners family demonstrate remarkable stability across challenging scenarios involving assumption violations and high censoring rates. We find that the choice of imputation method markedly impacts the performance of outcome imputation approaches, with \textcolor{red}{Margin} imputation consistently outperforming alternatives. 
    % Meanwhile, Causal Survival Forests shows stable and competitive performance without requiring extensive tuning of base model options, providing reliable "out-of-the-box" options for practitioners, though it struggles notably with positivity violations. \textcolor{red}{RED: not yet established; consider removing takeaways from experiments for space limit?}
\end{itemize}

While prior work has explored subsets of these design choices (e.g., \citet{ch04:cui2023estimating,meir2025heterogeneous}), \mbox{\textsc{SurvHTE-Bench}} is the first to systematically evaluate survival HTE methods under assumption violations, diverse survival models, and across synthetic, semi-synthetic, and real data.
We focus on binary treatments and static covariates with right-censored outcomes, as even this basic setting lacks a standardized benchmark. More complex extensions (time-varying treatments, longitudinal covariates, and instrumental variables) are beyond our present scope. 


%%%%%%%%%%%%%%%%%%%%%%%%%%%%%%%%%%%%%%%%%%%%%%

